\documentclass[../../main.tex]{subfiles}
\begin{document}

\begin{center}
  {\fontsize{48pt}{58pt}\selectfont \textbf{1974}\normalsize\textbf{年度}}
\end{center}

\begin{table}[htb]
  \centering
  \caption{出題テーマと難易度}
  \begin{tabular}{|c|p{9cm}|c|}
    \hline
    \textbf{問題番号} & \textbf{テーマ} & \textbf{難易度} \\
    \hline
    第1問 & 単位円上の3点と三角関数の和 & \\
    \hline
    第2問 & 半球形の容器から水をくみ出す速さ & \\
    \hline
    第3問 & 四次曲線のくぼみと極値の消滅 & \\
    \hline
    第4問 & 多項式とその導関数の整除関係 & \\
    \hline
    第5問 & 整数の平方と不等式 & \\
    \hline
    第6問 & 相似三角形の列の成立条件と無限和 & \\
    \hline
  \end{tabular}
\end{table}

\begin{table}[htb]
  \centering
  \caption{解答}
  \renewcommand{\arraystretch}{1.4}
  \begin{tabular}{|c|c|p{9cm}|}
    \hline
    \textbf{問題番号} & \textbf{小問} & \textbf{答え} \\
    \hline
    第1問 & -- & $\beta-\alpha=\gamma-\beta=\dfrac{2\pi}{3}$ \\
    \hline
    第2問 & (1) & $\pi\left(100vt-\dfrac{v^3t^3}{3}\right)\,\mathrm{cm^3}$ \\
    \cline{2-3}
    & (2) & $\pi v(100-v^2t^2)\,\mathrm{cm^3/秒}$ \\
    \hline
    第3問 & -- & 閾値は$q=p$ \\
    \hline
    第4問 & -- & $F(x)=\dfrac{x^{n+1}}{n+1}$，$g(x)=x^n$ \\
    \hline
    第5問 & (イ) & 条件を満たす整数の組は有限個（証明問題） \\
    \cline{2-3}
    & (ロ) & $m\ge9,a=8,b=9$では$n^2=4m+1$ \\
    \cline{2-3}
    & (ハ) & $(m,n)=(72,17)$ \\
    \hline
    第6問 & -- &  \\
    \hline
  \end{tabular}
\end{table}

\paragraph{試験形態}
6問．試験時間は要確認．

\end{document}
